\scanpage{91}\mspage{35}
% L1-2
qu'on peut regarder \unsure{donc|alors} \inserted{(via la construction $C$)} comme définissant un type d'homotopie relatif \uline{loc constant} sur $S$.

% L3-7
Donc pour élucider la situation i.e. les relations entre les 13 catégories localisées provenant des diagrammes de la page 8, il reste à regarder si les assertions hypothétiques que je viens de faire sont bien valables, savoir

% L8-9
\textcircled{\scriptsize 1\textsuperscript{o}} Le foncteur déduit de
\[
 \underline{\mathrm{Fibsc}}_0\,\deleted{\unsure{E}}S\longrightarrow\underline{\mathrm{Fib}}_0\,S
\]
% L10-15
par passage aux localisés est une équivalence de catégories (et on peut dans cette question remplacer $\underline{\mathrm{Fib}}_0\,S$ par n'importe quelle \inserted{autre} sous catégorie pleine de $\underline{\mathrm{Cat}}_{/S}$ la contenant -- cela revient au même), et

% L16-17
\textcircled{\scriptsize 2\textsuperscript{o}} Le foncteur déduit de
\[
 \underline{\mathrm{Fibsc}}_0\,\deleted{\unsure{E}}S\longrightarrow\underline{\mathrm{Fibsc}}\,S
\]
% L18-22
par localisation est pleinement fidèle. (Que son image \deleted{essentielle} soit exactement formée des objets loc. \unsure{constants} \unsure{du premier|des premiers} \unsure{localisé|localisés}, interprété comme $\mathrm{HOT}_W(S)$, est \deleted{\unsure{doute}} logique.)\unsure{.}

% L23
Il n'y a plus lieu de se soucier du côté
% archive sheet number: 90
